% 由 sympy_validation.py 生成；不要手工编辑。
\begin{longtable}{p{0.12\linewidth}p{0.20\linewidth}p{0.13\linewidth}p{0.45\linewidth}}
\toprule
编号 & 检查对象 & 状态 & 证据边界\\
\midrule
\endhead
SYM-01.01 & 量纲指数相加 & 通过 & 量纲一致只是否定错误的必要条件，不证明动力学公式。\\
SYM-01.02 & 对数换底 & 通过 & 不验证数据是否真的服从幂律或指数律。\\
SYM-01.03 & Euler 公式与单位模 & 通过 & 复相位的物理可观测性还需具体算符和对称性。\\
SYM-01.04 & 幂函数求导 & 通过 & 数值差分的舍入误差不在符号检查内。\\
SYM-01.05 & 微积分基本定理的多项式实例 & 通过 & 一般函数需满足可积性和绝对连续等分析条件。\\
SYM-02.01 & 正交变换保持长度 & 通过 & 有限维正交代理；Lorentz 变换使用不同度规。\\
SYM-02.02 & 二阶矩阵逆 & 通过 & 大稀疏矩阵的稳定求解和条件数需数值验证。\\
SYM-02.03 & 对称二阶矩阵谱分解 & 通过 & 一般厄米谱定理由线性代数证明；近简并数值稳定性另测。\\
SYM-02.04 & 二维 Levi--Civita 缩并 & 通过 & 高维缩并的符号取决于指标位置和度规约定。\\
SYM-02.05 & SU(2) 一参数子群 & 通过 & 仅验证 SU(2) 对角子群，不代替一般 SU(3) 群算法。\\
SYM-03.01 & N=4 离散 Fourier 正交性 & 通过 & 一般 N 的几何级数证明在课件中给出；FFT 轴与符号仍需实现测试。\\
SYM-03.02 & 高斯 Fourier 变换 & 通过 & 有限盒子、离散采样和混叠需数值收敛测试。\\
SYM-03.03 & 有限样本总体方差的仿射变换 & 通过 & 抽样方差、协方差估计和自相关另行处理。\\
SYM-03.04 & 中心差分 Taylor 展开 & 通过 & 一般光滑函数余项为 O(h^6)（除以 h 后对应式中 O(h^6) 之后）。\\
SYM-03.05 & 复合梯形法二次函数误差 & 通过 & 浮点舍入、一般函数高阶导数界和自适应策略不在此恒等式内。\\
SYM-04.01 & 1+1 维 Lorentz 间隔 & 通过 & 只验证标准 boost；一般四维协变由矩阵度规条件给出。\\
SYM-04.02 & 四动量质壳 & 通过 & 格点色散伪影不由连续代数式给出。\\
SYM-04.03 & 谐振子 Euler--Lagrange 方程 & 通过 & 只验证变分结果的代数形式；泛函分析条件在正文说明。\\
SYM-04.04 & 谐振子能量守恒 & 通过 & Noether 定理的一般场论证明不由单一模型替代。\\
SYM-04.05 & Abelian 场强反对称与规范不变 & 通过 & 非阿贝尔交换子项在 V07 单独验证。\\
SYM-05.01 & 高斯波函数归一化 & 通过 & 有限盒子和离散积分需独立收敛检查。\\
SYM-05.02 & 位置—动量对易子 & 通过 & 无界算符的定义域问题需泛函分析；离散导数有边界残差。\\
SYM-05.03 & 厄米二能级幺正演化 & 通过 & 时间依赖 Hamiltonian 需要时间有序指数。\\
SYM-05.04 & 二能级 Born 概率与期望 & 通过 & Born 规则是量子理论公设，SymPy 只核对其代数后果。\\
SYM-05.05 & 高斯最小不确定波包 & 通过 & 一般 Robertson 不确定关系的证明在课件正文；这里只验证饱和实例。\\
SYM-06.01 & Klein--Gordon 平面波色散 & 通过 & 自由场检查不证明相互作用 QCD 的谱。\\
SYM-06.02 & 截断 Fock 空间的产生算符 & 通过 & 有限维空间不可能在最高态满足完整 Heisenberg 代数；边界残差被显式排除。\\
SYM-06.03 & 高斯生成泛函 & 通过 & 验证有限维高斯模型及源导数；一般相互作用路径积分不作形式评价。\\
SYM-06.04 & 零均值高斯四点 Wick 配对 & 通过 & 多点自由场的三种配对结构由同一高斯定理给出；费米负号另需 Grassmann 代数。\\
SYM-06.05 & 欧氏谱的基态投影 & 通过 & 若基态重叠为零或有限时间边界显著，主导态需重新判定。\\
SYM-07.01 & 质子与中子电荷计数 & 通过 & 颜色单态和强相互作用动力学不由电荷计数证明。\\
SYM-07.02 & SU(3) 生成元恒等式 & 通过 & 用八个反厄米 Gell--Mann 生成元逐项验证李代数；约定转换在课件注明。\\
SYM-07.03 & U(1) 与非对易 SU(2) 协变导数 & 通过 & 精确验证非阿贝尔矩阵次序和非齐次项符号；四维 SU(3) 离散实现仍需逐点测试。\\
SYM-07.04 & 非阿贝尔场强与 QCD 拉氏量不变量 & 通过 & 矩阵代理验证交换子、自旋无关颜色双线性与迹型胶子项的不变性；完整 SU(3) 味结构、旋量指标和量子反常仍需场论推导。\\
SYM-07.05 & 一圈运行耦合方程 & 通过 & 靠近 Lambda 的强耦合和禁闭不能由一圈公式证明。\\
SYM-08.01 & 周期边界的动量量子化 & 通过 & 有限格点还需识别 k 模 N 和 Brillouin 区。\\
SYM-08.02 & 格点链接变量连续展开 & 通过 & 在可交换单色分量验证幺正性和 a^2 展开；非阿贝尔使用矩阵指数。\\
SYM-08.03 & U(1) plaquette 的离散 curl & 通过 & 非阿贝尔 Baker--Campbell--Hausdorff 展开和路径基点另行验证。\\
SYM-08.04 & Wilson plaquette 小场展开 & 通过 & SU(3) 迹归一化和格点求和系数依赖生成元约定。\\
SYM-08.05 & 连续与无限体积极限 & 通过 & 具体系数、幂次和有限体积函数须由作用量、算符与数据决定。\\
SYM-09.01 & Pauli 子代数的 Euclidean Clifford 关系 & 通过 & 四维 gamma 表示和项目轴顺序由共享约定测试保证。\\
SYM-09.02 & 朴素费米子倍增零点 & 通过 & 四维 16 重倍增由四方向笛卡尔组合得出。\\
SYM-09.03 & Wilson 项的物理支与倍增支 & 通过 & Wilson 项的手征破缺、临界质量和 Clover 改进需完整格点分析。\\
SYM-09.04 & 传播子线性系统的小矩阵闭合 & 通过 & 巨大稀疏 Dirac 算符的条件数、预条件与浮点停止准则需数值测试。\\
SYM-09.05 & Jacobi 到 Gaussian 涂抹极限 & 通过 & 自由周期小格点验证 Laplacian 模、极限与投影；非阿贝尔协变性另测。\\
SYM-10.01 & 两态 Boltzmann 归一化 & 通过 & Markov 链样本的无偏性还需平稳、遍历和热化条件。\\
SYM-10.02 & Metropolis--Hastings 两分支详细平衡 & 通过 & 遍历性和实际混合速度不由详细平衡单独保证。\\
SYM-10.03 & Markov 链自相关与有效样本数 & 通过 & 验证有限自回归代理的自相关和方差放大；真实链需窗口与误差分析。\\
SYM-10.04 & HMC 标量 leapfrog 与接受校正 & 通过 & 验证一维标量代理的可逆/保体积及 Hamiltonian 误差；SU(3) 力另测。\\
SYM-10.05 & 系综物理尺寸与 m\_pi L & 通过 & m\_pi L 大于某经验值不构成有限体积误差证明。\\
SYM-11.01 & 三夸克反对称颜色单态 & 通过 & 一般 SU(3) 结论依赖 epsilon U U U=det(U)epsilon；旋量/费米符号另测。\\
SYM-11.02 & 二点谱的基态投影 & 通过 & 真实重叠可有投影与归一化因子；有限时间边界须加入拟合。\\
SYM-11.03 & 单指数有效质量 & 通过 & 多态、噪声导致的非正关联函数与后向项需相关拟合。\\
SYM-11.04 & 三点中心插入的激发态抑制 & 通过 & A、B 的关系取决于算符与运动学，未强制对称；实际 K\_Pi 来自 Euclidean 自旋和与目标 Lorentz 投影，不能由标量代理猜测。\\
SYM-11.05 & 断连真空扣除的有限样本协方差结构 & 通过 & 验证 1/N 均值差式须乘 N/(N-1)；非零物理信号需要真实系综配对，N<2 必须由接口拒绝。\\
SYM-12.01 & 协方差的中心与原始矩形式 & 通过 & 无偏样本协方差改用 1/(N-1)；自相关需块处理。\\
SYM-12.02 & 样本均值的删除一 jackknife 方差 & 通过 & 非线性估计量关系只在大样本下一阶成立；课程算法仍重跑完整链。\\
SYM-12.03 & N=2 bootstrap 均值的精确枚举 & 通过 & 真实分析的块 bootstrap 和分位数覆盖需重复模拟验证。\\
SYM-12.04 & 相关 chi-square 与 nuisance profile & 通过 & 在有限维精确模型验证驻点、Woodbury 形式和 profile；协方差估计误差另测。\\
SYM-12.05 & 正能隙参数化 & 通过 & 参数有序不保证多态模型可由有限数据辨识。\\
SYM-13.01 & 质量可忽略部分子的 Bjorken x & 通过 & 靶质量、高 twist 和有限 Q^2 修正未由该运动学代理涵盖。\\
SYM-13.02 & 有限支持均匀分布的 Ioffe-time 变换 & 通过 & 光锥算符定义、Wilson 线和重整化由物理理论给出，非此积分证明。\\
SYM-13.03 & PDF Mellin 矩与局域算符系数 & 通过 & 有限代理验证矩关系和最低胶子动量矩；不存在由本记录支持的 n=1 规范不变胶子数算符，真实 mixing/renormalization 另验。\\
SYM-13.04 & Mellin 卷积代数 & 通过 & 验证可积测试函数的卷积和矩性质；DGLAP plus 分布与截断阶另有边界。\\
SYM-13.05 & 二维 TMD Fourier 归一化 & 通过 & 验证高斯/有限维代理的 bT--kT Fourier 关系；soft 与 rapidity 结构不由 Fourier 代数消除。\\
SYM-14.01 & quasi/pseudo Fourier 反演 & 通过 & 验证受控测试函数的正反变换；有限 z 重建的病态性需数值覆盖测试。\\
SYM-14.02 & LaMET 匹配与幂计数 & 通过 & 验证代理核的卷积、尺度和极限；修正系数一般依 x、bT、ell，且小 x/端点展开可能非一致，QCD 因子化定理本身不是 SymPy 结论。\\
SYM-14.03 & 两通道胶子—singlet 线性混合代理 & 通过 & 真实匹配核含分布、尺度和卷积；这里只检查通道线性代数。\\
SYM-14.04 & pseudo-ITD Fourier 关系 & 通过 & 验证有限测试分布的 Ioffe-time 变换和归一化；短距因子化域另判。\\
SYM-14.05 & Collins--Soper 演化积分 & 通过 & 只验证固定 bT 的演化积分；相同 z、不同 P\_z 不满足共同 nu，两点比值也不能辨认 K 与 1/P\_z^2 污染。\\
SYM-15.01 & clover 反厄米无迹投影代理 & 通过 & clover 的 8ga^2 系数和 O(a^2) 改进需小场格点展开验证。\\
SYM-15.02 & 双场强闭合颜色迹的端点规范不变 & 通过 & 一般 SU(3) 由相同矩阵消去和迹循环性证明；代数不区分 [-,-] 与其他 link class，过程标签和路径几何必须另测。\\
SYM-15.03 & 三段 staple 的端点位移 & 通过 & 端点代数验证 z=0,bT!=0 仍双局域；真正局域还要求 ell=0 和 Wilson 路径收缩，非零守卫与周期边界由实现检查。\\
SYM-15.04 & 二维横向张量的非极化、螺旋度与线偏振分解 & 通过 & 这里只验证横向张量代数；light-front 指标、i 因子、核子自旋和 f1g[-,-] 归一化必须由选定算符约定固定。\\
SYM-15.05 & 算符混合矩阵的正反映射 & 通过 & 真实 Z 是尺度、方案和几何相关的卷积/混合对象。\\
SYM-16.01 & 四维热核与半群 & 通过 & 验证归一化、卷积半群和扩散宽度代理；有限周期盒边界另测。\\
SYM-16.02 & Wilson 格点流的作用量单调性 & 通过 & 在显式格点代理验证梯度平方结构；生产积分器还需步长收敛。\\
SYM-16.03 & 流时窗口的无量纲尺度比 & 通过 & 不等式是否有重叠窗口取决于实际 a、几何和数据稳定性。\\
SYM-16.04 & flowed 传播子的热核抑制 & 通过 & 验证微扰线性化传播子的指数抑制；rapidity subtraction 明确不在该检查中。\\
SYM-16.05 & 局域小流时展开与非局域路径边界 & 通过 & 局域常数加 tau 模型不能验证非局域 TMD 的 flow-to-standard 转换；后者必须给出依赖 z、bT、ell、端点、cusp、mixing 与方案的 C\_Gamma 核，否则保持 finite-flow prototype。\\
SYM-17.01 & 流轨迹检查点的半群复用 & 通过 & 非线性群值 Wilson 流由数值积分器回归测试，不由单模代理替代。\\
SYM-17.02 & 算符几何网格计数 & 通过 & 路径缓存可减计算量但不能改变逻辑数据项数。\\
SYM-17.03 & boost 核子连续色散 & 通过 & 有限 a 的格点色散和信噪比须由多动量数据检查。\\
SYM-17.04 & 构型级无偏断连协方差 & 通过 & SymPy 验证 1/(N-1) 中心化形式与 N/(N-1) 均值差式等价；N<2 必须由接口显式失败，构型错配仍需故障注入测试。\\
SYM-17.05 & 双态领先项的 summation 几何级数 & 通过 & 真实多态、相关噪声和有限 T 需共享 C2 参数的联合拟合。\\
SYM-18.01 & Wilson 线线性反项指数化 & 通过 & 验证长度可加自能的指数形式；端点、cusp、flow-time 与方案依赖另列。\\
SYM-18.02 & 已声明方案中的对称 quasi-soft 平方根分配 & 通过 & 只核对特定对称约定的代数。任意 Euclidean 真空 staple 不自动是标准 soft；缺 C\_qsoft\_to\_S、方向、regulator 或 zero-bin 时必须保持 prototype。\\
SYM-18.03 & 比值型非微扰重整化 & 通过 & 验证共同乘法因子的消去与参考条件；参考态 IR 系统学不自动消失。\\
SYM-18.04 & 直线算符的 hybrid 重整化连续拼接 & 通过 & 该检查严格限定为 straight-line hybrid。finite-staple 必须以完整路径周长差 L\_Gamma-L\_Gamma,s 重新推导，并另处理端点、cusp 与 mixing；不得由本记录升级为已重整化 TMD。\\
SYM-18.05 & quasi-TMD 匹配与 CS 核 & 通过 & 底层有限模型只验证 leading-power 两点比值代数，不是生产充分条件；完整 f1g[-,-] 胶子-singlet 匹配、共同运动学、三动量联合拟合与真实数据闭合仍须另验。\\
SYM-19.01 & a^2 连续外推截距 & 通过 & 实际幂次、尺度相关和 a^2/tau 交叉项由作用量与数据决定。\\
SYM-19.02 & 大动量倒幂外推 & 通过 & 匹配对数须先统一；d\_n 一般依赖 x、b\_T、ell，且小 x 与端点的倒幂展开可能非一致；(aP)^2 与 1/P^2 需交叉数据区分。\\
SYM-19.03 & 连续后零流时的有序极限 & 通过 & safe\_order 只属于局域或已完成 C\_Gamma 转换的对象；未消去 L\_Gamma/sqrt(8tau)、ln tau、cusp 与 mixing 时，非局域 TMD 必须触发停止门。\\
SYM-19.04 & 有限 staple 指数饱和 & 通过 & 增大 ell 只测试有限长度饱和，不会把空间型方向变为光状或替代 rapidity 方案；幂尾、周期镜像和 soft 几何效应必须与替代模型比较。\\
SYM-19.05 & 相关误差来源的总方差 & 通过 & 系统变体间相关性须由联合重采样或明确模型估计。\\
SYM-20.01 & 命名轴笛卡尔积的元素数 & 通过 & 分块、压缩和稀疏布局只改变物理存储，不改变逻辑契约。\\
SYM-20.02 & 内容寻址依赖的符号结构 & 通过 & 碰撞抗性是密码学性质，不由 SymPy 证明；实际实现还需规范化字节编码和原子写入。\\
SYM-20.03 & 标准正态两倍标准差覆盖率 & 通过 & 真实覆盖率必须由重复合成实验测量，单次 |pull|<2 不证明正确。\\
SYM-20.04 & Amdahl 强标度 & 通过 & 真实性能必须 profile；模型不是 PyQCD 实测。\\
SYM-20.05 & 四类调节量的联合理想极限 & 通过 & 真实极限含交叉项、对数、混合和相关数据；该检查只验证终章流程骨架。\\
SYM-21.01 & 群作用、子群、陪集与 Lagrange 定理 & 通过 & 矩阵实例精确验证 e·x=x 与 (gh)·x=g·(h·x)；一般群作用、轨道稳定子和 Lagrange 定理由集合双射证明。\\
SYM-21.02 & 表示、不可约表示与 Schur 引理三情形 & 通过 & 有限实例分别核对 Schur 引理的不等价、同一表示与等价副本三种情形；一般证明仍依不可约表示的核与像。\\
SYM-21.03 & 特征标正交与表示分解 & 通过 & 验证 C3 精确实例；一般群需完整共轭类和特征标表。\\
SYM-21.04 & 立方旋转群 O、O\_h、宇称与自旋下沉 & 通过 & 精确完成 J=2↓O=E⊕T2 的特征标内积；实际格点自旋指认仍需宇称、能级简并、连续外推与重叠。\\
SYM-21.05 & 有限群投影算符 & 通过 & 验证 C2 投影实例；一般公式依矩阵元/特征标正交。\\
SYM-22.01 & 微分几何桥梁：流形、切空间与指数映射 & 通过 & 矩阵指数实例验证局部群结构；一般李群的全局指数映射性质未证明。\\
SYM-22.02 & 厄米生成元、结构常数、Jacobi 恒等式与 BCH & 通过 & 验证 BCH 在中心交换子实例上精确截断；一般级数收敛域另论。\\
SYM-22.03 & SU(2) 角动量表示与 Casimir & 通过 & 验证一个不可约表示；半整数表示与群全局性质需另查。\\
SYM-22.04 & SU(3) 根、权、Dynkin 标号、维数与 Casimir & 通过 & 精确验证 A2 简单根、基本权、维数与二次 Casimir 公式的 3 和 8 表示；权重 multiplicity 与一般张量积分解仍需表示论算法。\\
SYM-22.05 & 测度桥梁、Haar 积分与矩阵元正交 & 通过 & SymPy 只解归一化常数；Haar 存在唯一性和不变量张量由群论证明。\\
SYM-23.01 & 一代标准模型场表、手征性与电荷 & 通过 & 只核对电荷归一化；规范反常与完整粒子谱另验。\\
SYM-23.02 & 夸克模型、flavor 多重态与强子分类 & 通过 & 维数相等是必要检查，不单独证明不可约分解。\\
SYM-23.03 & 角动量耦合、宇称、荷共轭与 JPC & 通过 & 不验证 flavor 相位和非自共轭强子的 C 标签。\\
SYM-23.04 & 两体相空间、散射截面与衰变宽度 & 通过 & 验证通量/两体相空间给出的 dσ/dΩ 与积分；自旋颜色平均、相同粒子因子、质量阈值和真实振幅需逐过程处理。\\
SYM-23.05 & 显式卷积、硬散射因子化与部分子模型 & 通过 & 只验证一个显式卷积和高尺度幂修正；QCD 因子化的区域分离、Glauber 条件与核的微扰计算不是 CAS 命题。\\
SYM-24.01 & 正则量子化、模展开与真空 & 通过 & 验证单自由度 [x,p]=i；场的分布 δ 与真空发散未由 SymPy 处理。\\
SYM-24.02 & 复分析与分布桥梁：传播子、留数和 iε & 通过 & 不验证 Fourier 轮廓、因果边界和分布意义。\\
SYM-24.03 & 相互作用图、生成泛函与 Feynman 规则 & 通过 & 只验证四点 Wick 配对数 3；真实时空图、对称因子和正规化另算。\\
SYM-24.04 & 固定归一化的 Minkowski LSZ 与 Euclidean 边界 & 通过 & 验证 LSZ 截肢的极点代数；渐近态存在和多粒子支切另论。\\
SYM-24.05 & λφ⁴ 一圈：正则化、极点、反项与 beta & 通过 & 完整验证一圈 bubble 的极点、delta\_lambda 与 beta\_lambda=3lambda^2/(16pi^2)；有限部分、二圈、阈值与非微扰 triviality 不在本检查中。\\
SYM-25.01 & Hamilton 桥梁：相空间、约束与两个物理偏振 & 通过 & 只验证常场散度为零；非 Abelian 约束代数与物理 Hilbert 空间另论。\\
SYM-25.02 & Grassmann 桥梁与 Faddeev--Popov 行列式 & 通过 & 不处理无限维行列式、Gribov copies 或正则化。\\
SYM-25.03 & 完整 BRST、分次 Leibniz 与规范固定作用量 & 通过 & 有限矩阵实现验证两个 Grassmann 生成元和分次符号，并核对 s[bar\_c(G+xi B/2)]；完整时空作用量、测度与 Gribov 问题仍需场论处理。\\
SYM-25.04 & Z1、Z2、Ward 与 Slavnov--Taylor 恒等式 & 通过 & 只验证树级纵向恒等式；圈级 Slavnov--Taylor 含鬼核与重整化。\\
SYM-25.05 & 轴反常统一约定与一代标准模型反常抵消 & 通过 & 精确核对四类局域规范反常及课程轴 Ward 恒等式系数；全局反常、正则化推导和非微扰拓扑涨落不由电荷求和证明。\\
SYM-26.01 & 自发对称破缺与 Goldstone 定理 & 通过 & 有限体积自发破缺与量子修正未包含。\\
SYM-26.02 & Higgs 机制与规范玻色子质量 & 通过 & 不证明规范不变可观测谱或非 Abelian 混合。\\
SYM-26.03 & 有效场论、OPE、matching 与系统幂计数 & 通过 & 只验证 Λ^{4-d} 配平；Wilson 系数大小与 EFT 幂计数需物理输入。\\
SYM-26.04 & 拓扑桥梁：S³∞→SU(N)、π3=Z 与 instanton & 通过 & 精确验证紧化空间的显式映射落在 SU(2)；pi\_3(SU(N))=Z、绕数积分整数性与 instanton 作用量界仍由拓扑和场论证明。\\
SYM-26.05 & Wilson 圈、Polyakov loop 中心对称与禁闭边界 & 通过 & 验证 Wilson 面积律代理及 L\_P->zL\_P、局域 plaquette 不变；动力学基本表示夸克显式破坏中心对称，弦断裂和 crossover 需有限温度数据。\\
SYM-27.01 & 转移矩阵与二点函数谱分解 & 通过 & 只验证两态比值；真实谱权重、边界和噪声需数据拟合。\\
SYM-27.02 & 有效质量、周期边界与多态拟合 & 通过 & arccosh 主支和含噪非线性偏差需数值处理。\\
SYM-27.03 & 相关矩阵与广义本征值问题 & 通过 & 有限算符基修正、条件化和近简并态追踪不由对角代理覆盖。\\
SYM-27.04 & LapH distillation、perambulator 与算符因子化 & 通过 & 不验证协变 Laplacian、perambulator 求逆或规范协变性。\\
SYM-27.05 & 立方群下沉、算符重叠与连续自旋指认 & 通过 & 立方群系数与连续自旋的物理指认仍需特征标和谱数据。\\
SYM-28.01 & 两粒子运动学与非相互作用能级 & 通过 & 只验证连续色散运动学；实际格点应使用实测/格点色散。\\
SYM-28.02 & 弹性 Lüscher 量化条件与相移 & 通过 & 不计算 Z00；一般 moving-frame/分波混合需矩阵条件。\\
SYM-28.03 & 有效程展开、散射长度与束缚态 & 通过 & 近阈值参数化的收敛和 Riemann 片层需物理论证。\\
SYM-28.04 & moving frame、little group 与分波混合 & 通过 & 真实盒矩阵 B、little group 和 lmax 截断需数值实现。\\
SYM-28.05 & 耦合道振幅、解析延拓与共振极点 & 通过 & 只验证极点参数定义；耦合道幺正、解析延拓和留数需完整振幅。\\
SYM-29.01 & Poisson 求和与有限盒修正的来源 & 通过 & 验证 Poisson 公式中的单个 Fourier 像；无限格点求和收敛另论。\\
SYM-29.02 & 质量隙、mπL 与单强子指数修正 & 通过 & 观测量系数、多绕行和更重态未验证。\\
SYM-29.03 & 两粒子在壳传播与幂律体积效应 & 通过 & a0/L 高阶、有效程和完整 Lüscher 条件未包含。\\
SYM-29.04 & 扭曲边界条件与连续动量调谐 & 通过 & 不验证 partial twisting 的海夸克和 annihilation 有限体积效应。\\
SYM-29.05 & QED 有限体积、零模与方案依赖 & 通过 & 不固定 QED\_L/QED\_TL 的普适系数或结构项。\\
SYM-30.01 & 热迹、Euclidean 时间圆与温度设定 & 通过 & 不验证路径积分测度和费米反周期迹符号。\\
SYM-30.02 & Polyakov loop、中心对称与自由能 & 通过 & 静态自由能重整化和动力学夸克显式破缺未验证。\\
SYM-30.03 & 手征凝聚、susceptibility 与 crossover & 通过 & 真实 stochastic trace、加性/乘性重整化和临界缩放未处理。\\
SYM-30.04 & 状态方程、trace anomaly 与积分法 & 通过 & 格点 beta functions、零温扣除和积分常数需实测。\\
SYM-30.05 & 热谱函数、输运系数与病态反演 & 通过 & 不证明连续谱唯一性；先验偏差必须用 mock-data 覆盖。\\
SYM-31.01 & staggered 费米子、taste 与 rooting & 通过 & rooting 的局域性和正确 universality 不能由质量极限单独证明。\\
SYM-31.02 & domain-wall 费米子与剩余质量 & 通过 & dislocation 尾、M5 调谐和 Ward identity 需真实数据。\\
SYM-31.03 & overlap 算子与 Ginsparg--Wilson 关系 & 通过 & 不验证 overlap 局域性、低模近似误差或 index theorem。\\
SYM-31.04 & twisted-mass Wilson 与最大扭转自动改进 & 通过 & 不验证 PCAC 临界质量调谐和 flavor/parity O(a²) 劈裂。\\
SYM-31.05 & mixed-action、partial quenching 与统一连续极限 & 通过 & partial-quenching double poles 与 EFT 高阶未验证。\\
SYM-32.01 & Symanzik 改进规范作用量 & 通过 & 不验证 action 的 BCH 全展开、圈级改进或实际旋转破缺。\\
SYM-32.02 & 各向异性格点与色散调谐 & 通过 & 重整化各向异性必须由实测色散调谐。\\
SYM-32.03 & 拓扑冻结、自相关与算法临界变慢 & 通过 & 真实 τint 的窗口选择、慢尾和拓扑遍历需时间序列诊断。\\
SYM-32.04 & 开放时间边界与边界污染 & 通过 & 有限 T 的 min 距离、边界反项和 bulk 平台需数据验证。\\
SYM-32.05 & 尺度设定：r0、t0、w0 与误差传播 & 通过 & 不确定度相关传播和不同尺度间的连续一致性需重采样。\\
SYM-33.01 & 重整化矩阵、算符混合与 MSbar & 通过 & 实际算符基、power mixing 和 perturbative truncation 未验证。\\
SYM-33.02 & RI/MOM：固定规范下的非微扰重整化 & 通过 & Landau 规范、Goldstone pole、H(4) 和窗口问题需格点数据。\\
SYM-33.03 & RI/SMOM：非例外动量与 Ward 恒等式 & 通过 & 只验证非例外运动学；projector 与 MSbar conversion 需具体方案。\\
SYM-33.04 & Schrödinger functional 与算符 step scaling & 通过 & 每步的调谐、相关误差与 a/L 连续外推未验证。\\
SYM-33.05 & gradient-flow coupling 与有限体积 beta function & 通过 & 不验证 flow coupling 的 tree normalization、scheme 或连续外推。\\
SYM-34.01 & Wilson-line 线性发散与乘法结构 & 通过 & 端点/cusp 对数、mixing 与 flow scheme 仍需独立处理。\\
SYM-34.02 & 辅助场方法：Wilson line 化为静态传播子 & 通过 & 不验证路径有序、表示、cusp 或端点 current mixing。\\
SYM-34.03 & ratio 与自重整化：从同源数据消去 UV 因子 & 通过 & 不同长度/几何的 Z 比、IR 参考依赖和绝对匹配未验证。\\
SYM-34.04 & hybrid 方案、切换尺度与连续拼接 & 通过 & 指数符号、zs 窗口、导数平滑和 Fourier 系统学需方案扫描。\\
SYM-34.05 & 方案标记的 quasi-TMD、soft、CS kernel 与 flow 转换 & 通过 & 不证明 soft 几何匹配、rapidity 因子化、flow conversion 或胶子 matching。\\
SYM-35.01 & 从 $f\_1^{g[-,-]}$ estimand 到四级证据状态、DAG 与停止门 & 通过 & 实际 estimand 还需 scheme、soft、匹配、交叉项和可识别性。\\
SYM-35.02 & flow window 与 a、τ、L、Pz、ell 的分阶段层级极限 & 通过 & 精确验证协方差两极端、有序极限和三类安全距离都能主导 flow window；不证明非局域 C\_Gamma 的路径、端点、cusp、mixing 或方案转换，缺核时仍须停止在 finite-flow prototype。\\
SYM-35.03 & 有限 Ioffe-time 数据的 Fourier 逆问题 & 通过 & 物理约束、分辨核、νmax 和正则化覆盖需 mock-data 实测。\\
SYM-35.04 & 可复现管线、数据谱系与审计 & 通过 & 不证明密码学抗碰撞、规范化编码或并行 bitwise reproducibility。\\
SYM-35.05 & 最终资格考核：按四级证据独立实现 $f\_1^{g[-,-]}$ 链 & 通过 & 只证明均值乘积不能恢复 connected covariance；完整生产链仍需真实数据全门验证。\\
\bottomrule
\end{longtable}
